\documentclass[10pt,a4paper,twocolumn]{article}

\usepackage[top=1.8cm, bottom=1.8cm, left=1.5cm, right=1.5cm, columnsep=0.6cm]{geometry}
\usepackage{times}
\usepackage{microtype}
\usepackage{titlesec}
\usepackage{parskip}
\usepackage{booktabs}
\usepackage{enumitem}
\usepackage{xcolor}
\usepackage{hyperref}
\usepackage{fancyhdr}
\usepackage{lastpage}
\usepackage{amsmath}
\usepackage{setspace}

% Colors
\definecolor{navyblue}{HTML}{1a1a2e}
\definecolor{darkblue}{HTML}{16213e}
\definecolor{accentblue}{HTML}{0f3460}

% Section formatting
\titleformat{\section}{\normalfont\bfseries\fontsize{9.5}{11}\selectfont\color{darkblue}}{\thesection.}{0.4em}{}[\vspace{-4pt}\textcolor{darkblue}{\rule{\linewidth}{0.5pt}}]
\titleformat{\subsection}{\normalfont\bfseries\fontsize{9}{10.5}\selectfont\itshape}{\thesubsection}{0.4em}{}
\titlespacing{\section}{0pt}{8pt}{4pt}
\titlespacing{\subsection}{0pt}{5pt}{2pt}

% Paragraph spacing
\setlength{\parskip}{3pt}
\setlength{\parindent}{0pt}

% Header/footer
\pagestyle{fancy}
\fancyhf{}
\renewcommand{\headrulewidth}{0.4pt}
\fancyhead[L]{\fontsize{7}{8}\selectfont\textcolor{darkblue}{\textit{UAV-Based Post-Disaster Debris Detection --- Project Proposal}}}
\fancyhead[R]{\fontsize{7}{8}\selectfont\textcolor{darkblue}{MSDSM Batch 4 $\cdot$ 2404107067}}
\fancyfoot[C]{\fontsize{7}{8}\selectfont\thepage\ of \pageref{LastPage}}

% Hyperlink style
\hypersetup{colorlinks=true, urlcolor=accentblue, linkcolor=accentblue, breaklinks=true}

% Tighter itemize
\setlist[itemize]{noitemsep, topsep=2pt, leftmargin=*}

\begin{document}

% ─── Title block ────────────────────────────────────────────────────────────
\twocolumn[{%
\begin{center}
  {\fontsize{14}{17}\bfseries\color{navyblue}
    High-Resolution UAV-Based Post-Disaster Debris Detection:\\[2pt]
    Fine-Tuning Foundation Models on Multi-Source Aerial Imagery}\\[6pt]
  {\fontsize{9}{11}\selectfont\itshape\color{gray}
    Utkarsh Kaushik \quad $\cdot$ \quad MSDSM Batch 4 \quad $\cdot$ \quad Roll No.\ 2404107067 \quad $\cdot$ \quad 27 February 2026}\\[2pt]
  {\fontsize{8.5}{10}\selectfont\textcolor{gray}{Project Proposal --- Humanitarian AI for Disaster Response}}
\end{center}
\vspace{4pt}
\textcolor{darkblue}{\rule{\linewidth}{1pt}}
\vspace{6pt}
}]

% ─── 1. Introduction ────────────────────────────────────────────────────────
\section{Introduction}

Natural disasters cause devastating infrastructure damage, with hurricanes alone
inflicting over \textbf{\$50 billion} in annual losses in the United States.
Rapid, accurate debris assessment is critical for emergency response coordination,
yet traditional manual surveys are dangerous, time-consuming, and unscalable.
Unmanned Aerial Vehicles (UAVs) have emerged as essential tools for post-disaster
reconnaissance, generating high-resolution imagery that exceeds satellite capabilities
in spatial resolution and temporal flexibility.

This project addresses \textbf{automated debris detection and segmentation} from UAV
imagery using open-vocabulary vision-language models. Current state-of-the-art
models such as Florence-2 and SAM2 exhibit limited performance on disaster-specific
scenarios due to domain shift from natural imagery to debris-filled, structurally
compromised environments. We propose domain-specific fine-tuning of these foundation
models on curated post-disaster UAV datasets to enable: (1)~\textbf{text-guided
debris queries} without predefined categories, (2)~\textbf{pixel-accurate
segmentation} for damage quantification, and (3)~\textbf{rapid deployment} for
emergency response teams.

% ─── 2. Methodology ─────────────────────────────────────────────────────────
\section{Methodology}

The proposed architecture implements a \textbf{cascaded fine-tuning pipeline}
leveraging complementary strengths of vision-language and segmentation foundation
models.

\subsection{Florence-2 Fine-Tuning for Debris Detection}

Florence-2 (Microsoft, 2024) serves as the open-vocabulary detection backbone.
We employ \textbf{Low-Rank Adaptation (LoRA)} with rank $r{=}16$ and
$\alpha{=}32$ to efficiently adapt the model to disaster-specific vocabulary.
The training objective optimises detection of debris categories---including
damaged buildings, flooded regions, scattered rubble, and vegetation
debris---via task-specific prompting with
\texttt{\small<OPEN\_VOCABULARY\_DETECTION>} tokens.
Domain adaptation addresses irregular debris patterns, occlusion from structural
collapse, and heterogeneous lighting conditions in UAV-captured imagery.

\subsection{SAM-2 Fine-Tuning for Precise Segmentation}

SAM-2 (Meta AI, 2024) provides pixel-accurate mask generation guided by bounding
boxes from Florence-2. We adopt a \textbf{partial fine-tuning strategy}, freezing
the Hiera image encoder while training the prompt encoder and mask decoder on
aerial disaster imagery. This preserves general visual features while adapting
mask prediction to debris boundaries. The combined loss integrates Dice loss for
mask overlap and binary cross-entropy for boundary refinement.

\subsection{Integrated Inference Pipeline}

The deployed system processes UAV imagery through: (1)~Florence-2 open-vocabulary
detection with domain-adapted prompts, (2)~coordinate transformation for SAM-2
compatibility, (3)~mask generation with multi-scale refinement, and (4)~priority-based
filtering using a debris taxonomy (critical infrastructure damage, blocked access
routes, hazardous materials). Outputs include GIS-compatible GeoJSON, structured
JSON reports for emergency management systems, and visual priority maps.

% ─── 3. Datasets ────────────────────────────────────────────────────────────
\section{Datasets}

The proposed technique is evaluated on three complementary high-resolution UAV
datasets curated for post-disaster analysis.

\subsection{RescueNet}
\textit{Source:} Alam et al., IEEE TGRS, 2022.
\textit{Scale:} 4,494 images at 0.5--2\,cm/px GSD from Hurricane Michael (FL, 2018).
Pixel-level semantic segmentation with 8 classes: Background, Water,
Building-No-Damage, Building-Damaged, Vegetation, Road-No-Damage, Road-Damaged,
and Vehicle. Captures structural collapse, vegetation blockage, and road
obstruction from a Category~5 hurricane. Official train/val/test splits enable
standardised evaluation.

\subsection{DesignSafe-CI PRJ-6029}
\textit{Source:} NSF NHERI DesignSafe-CI Repository.
DOI: \href{https://doi.org/10.17603/ds2-jvps-2n95}{10.17603/ds2-jvps-2n95}.
Multi-hazard UAV and ground-based imagery from earthquakes, hurricanes, and floods
with georeferenced damage observations. Enables cross-disaster generalisation
testing; companion datasets support structural damage assessment and multi-modal
fusion research.

\subsection{MSNet Dataset}
\textit{Source:} Xie \& Xiong, Remote Sensing, 2023.
8,700+ images across 7 disaster events at 0.3--5\,cm/px.
Multi-scale damage classification (no damage, minor, major, destroyed) with
oriented bounding boxes and instance segmentation. Covers earthquakes (Chile,
Ecuador), hurricanes (Michael, Laura), and tsunami events, enabling evaluation
across varying UAV altitudes and sensor configurations.

\smallskip
\textbf{Evaluation Protocol.}~We adopt mean Intersection-over-Union (mIoU) for
segmentation, F1-score for debris detection, and Average Precision (AP) at IoU
$\in[0.5{:}0.95]$. Cross-dataset validation trains on RescueNet and tests on
DesignSafe-CI and MSNet holdout regions.

% ─── 4. Expected Outcomes ───────────────────────────────────────────────────
\section{Expected Outcomes}

This project anticipates four primary contributions to humanitarian computing:

\begin{itemize}
  \item \textbf{Domain-Adapted Foundation Models:} Publicly released Florence-2
    and SAM-2 checkpoints fine-tuned on post-disaster UAV imagery, enabling
    zero-shot deployment for resource-constrained emergency response organisations.
  \item \textbf{Cross-Dataset Benchmark:} Standardised evaluation protocol and
    baseline results across RescueNet, DesignSafe-CI, and MSNet, establishing
    reference metrics for high-resolution aerial debris detection.
  \item \textbf{Open-Source Deployment Framework:} Gradio-based web interface and
    Python API with export formats compatible with FEMA debris management systems
    and common GIS platforms.
  \item \textbf{Transfer Learning Analysis:} Empirical study of foundation model
    transfer across disaster types (hurricane vs.\ earthquake), informing data
    collection priorities for future dataset construction.
\end{itemize}

% ─── 5. References ──────────────────────────────────────────────────────────
\section{References}

\fontsize{7.5}{9.5}\selectfont
\setlength{\parskip}{2pt}

\noindent[1] M.\,S.\ Alam et al., ``RescueNet: A High-Resolution Post-Disaster UAV Dataset
for Semantic Segmentation,'' \textit{IEEE Trans.\ Geosci.\ Remote Sens.}, vol.~60, 2022.

\noindent[2] K.\ Amini et al., ``Hurricane-Induced Debris Segmentation Dataset Using Aerial
Imagery'' [v2], DesignSafe-CI, 2025. \href{https://doi.org/10.17603/ds2-jvps-2n95}{doi:10.17603/ds2-jvps-2n95}

\noindent[3] X.\ Zhu, J.\ Liang, and A.\ Hauptmann, ``MSNet: A Multilevel Instance
Segmentation Network for Natural Disaster Damage Assessment in Aerial Videos,''
\textit{WACV}, pp.~2023--2032, 2021.

\noindent[4] W.\ Xiao et al., ``Florence-2: Advancing a Unified Representation for a Variety
of Vision Tasks,'' \textit{CVPR}, 2024.

\noindent[5] N.\ Ravi et al., ``SAM~2: Segment Anything in Images and Videos,''
\textit{arXiv:2408.00714}, 2024.

\noindent[6] E.\,J.\ Hu et al., ``LoRA: Low-Rank Adaptation of Large Language Models,''
\textit{ICLR}, 2022.

\end{document}
